\documentclass[11pt,a4paper]{article}

\usepackage[T1]{fontenc}
\usepackage[margin=2cm]{geometry}
\usepackage{enumitem}

\setlength{\parskip}{0.8em}
\setlength{\parindent}{0pt}

\begin{document}

\begin{center}
    {\LARGE\textbf{Year 12 to Year 13 Transition Work}}\\[0.5em]
    {\Large Nuclear Power Research Guide}
\end{center}

The aim of this activity is to develop an understanding of how a modern nuclear power station operates. While watching the recommended video, produce your own notes covering the topics below.

You should not simply copy information from the video. Instead, aim to summarise the key ideas in your own words.

The online quiz will assess your understanding of the following topics and subtopics.  Make notes accordingly.

\bigskip

\section*{1. Nuclear Fuel}

\begin{itemize}[leftmargin=*]
    \item The fuel used in a modern nuclear reactor.
    \item The structure of a fuel assembly.
    \item Fuel rods and uranium fuel pellets.
    \item Why nuclear fuel is considered highly energy dense.
\end{itemize}

\section*{2. Nuclear Fission}

\begin{itemize}[leftmargin=*]
    \item What happens during nuclear fission.
    \item The role of neutrons.
    \item How energy is released.
    \item How a nuclear chain reaction is maintained.
\end{itemize}

\section*{3. The Reactor}

\begin{itemize}[leftmargin=*]
    \item The purpose of the reactor core.
    \item The function of the coolant.
    \item Why pressurised water reactors use separate primary and secondary water systems.
    \item The engineered barriers that prevent radioactive material escaping.
\end{itemize}

\section*{4. Heat Transfer}

\begin{itemize}[leftmargin=*]
    \item How heat is transferred away from the reactor core.
    \item The purpose of the steam generator.
    \item Why the reactor must continue to be cooled after shutdown.
\end{itemize}

\section*{5. Electricity Generation}

\begin{itemize}[leftmargin=*]
    \item How steam is produced.
    \item The purpose of the turbine.
    \item The purpose of the electrical generator.
    \item How the steam is recycled after passing through the turbine.
\end{itemize}

\section*{6. Radiation}

\begin{itemize}[leftmargin=*]
    \item The difference between fresh fuel and spent fuel.
    \item Why spent fuel remains radioactive.
    \item How radiation can affect living tissue.
    \item Why iodine-131 is considered particularly hazardous after a nuclear accident.
\end{itemize}

\section*{7. Radioactive Waste}

\begin{itemize}[leftmargin=*]
    \item What happens to spent nuclear fuel.
    \item How high-level radioactive waste is stored.
    \item Why the total volume of nuclear waste is relatively small.
\end{itemize}

\section*{8. Reactor Safety}

\begin{itemize}[leftmargin=*]
    \item The purpose of the containment building.
    \item The meaning of ``defence in depth''.
    \item Why refuelling operations require strict procedures.
\end{itemize}

\section*{9. The Nuclear Fuel Cycle}

\begin{itemize}[leftmargin=*]
    \item How fresh fuel arrives at a power station.
    \item How fresh fuel is prepared before entering the reactor.
    \item What happens during a routine refuelling outage.
\end{itemize}

\section*{10. Nuclear Power in Context}

\begin{itemize}[leftmargin=*]
    \item Why nuclear power is considered an efficient method of generating electricity.
    \item Advantages of modern reactor designs.
    \item The overall process by which a nuclear power station converts nuclear energy into electricity.
\end{itemize}

\end{document}